% Candidate complete continuation of Part B.  Not yet promoted to checked paper.
% Internal proof audit requested from the root agent, 2026-11-20 00:52 UTC.

\subsubsection*{Candidate B.2d: simultaneous embedding into one Jacobian}
Under the hypotheses and choices of B.2a, write $J=\operatorname{Pic}^0_{C/E}$.
For every proper normal modification $W\to M$, the unique lift
$C\to W_E$ induces a closed immersion
\[
 \operatorname{Pic}^{\tau}_{W/k}\times_kE\hookrightarrow J.
 \tag{B.2.3}
\]
These immersions commute with pullback between modifications.
\begin{proof}
The curve lies in the isomorphism locus of $W\to M$, and this determines
its lift uniquely.  Restriction takes torsion-equivalent line bundles to
degree-zero line bundles, so the displayed map is defined.  Its source
is a proper commutative group of finite type.  Its kernel has zero Lie
algebra by B.2b.  A finite type group scheme over the algebraically closed
field $E$ with zero Lie algebra is reduced of dimension zero: at the
identity its maximal ideal satisfies $\mathfrak m/\mathfrak m^2=0$,
hence $\mathfrak m=0$ by Nakayama's lemma, and translation treats every
closed point.  The kernel is proper, so it is finite \'etale.

If this kernel were nontrivial, its group of $E$-points would have a
nonzero element of prime order $\ell$.  Write
$P=\operatorname{Pic}^{\tau}_{W/k}$.  The group scheme $P[\ell]$ is
finite, even when $\ell=p$: on every component of $P_{\mathrm{red}}$,
multiplication by $\ell$ has finite fibres, because its identity component
is an abelian variety.  Thus $[\ell]:P\to P$ is proper and quasi-finite,
hence finite.  Every $E$-point of the finite $k$-scheme $P[\ell]$ is a
$k$-point, since $k$ is algebraically closed.  Consequently our element
is represented by a line bundle $L$ on $W$ of exact order $\ell$, whose
restriction to $C$ is trivial.

Suppose first that $\ell\ne p$.  The Kummer sequence, together with
$\Gamma(W,\mathcal O_W)^*=k^*$, gives a nontrivial
$\mu_\ell$-torsor $Q\to W$.  It is connected, since a torsor under
a cyclic group of prime order is either connected or trivial.  It is
normal and therefore integral.  Its normalization over $M$ in $k(Q)$
is finite and integral.  By B.2a its inverse image over $C$ is
geometrically irreducible, hence the \'etale cover $Q_E|_C$ is connected.
But $L|_C$ is trivial and $E^*$ is $\ell$-divisible, so this cover is
trivial.  This is a contradiction.

Now let $\ell=p$.  The Kummer sequence gives a nontrivial
$\mu_p$-torsor on $W$.  Restriction to the generic point is injective
for finite torsors on a normal integral scheme, by the normal full
faithfulness used in B.1b.  Its generic class is therefore represented
by $a\in k(W)^*$ with $a\notin k(W)^{*p}$.  The logarithmic differential
$da/a$ is nonzero, since the kernel of $d:k(W)\to
\Omega^1_{k(W)/k}$ is $k(W)^p$.  Local equations for the torsor are
$t^p=u$ with $u$ a unit; changing a trivialization multiplies $u$ by a
$p$-th power.  Hence $da/a$ is regular on $W_{\mathrm{reg}}$.
On the common open of $M$ and $W$ it extends, by Hartogs on the smooth
scheme $U$, to an element of $V=H^0(U,\Omega^1_{U/k})$.
The triviality of $L|_C$, and the surjectivity of $E^*\xrightarrow{p}E^*$,
make the restricted $\mu_p$-torsor trivial.  Thus $da/a$ restricts to
zero on $C$, contradicting B.2a(1).

The kernel is therefore trivial.  A homomorphism of group schemes with
trivial scheme-theoretic kernel is a monomorphism, and a proper
monomorphism is a closed immersion.  Compatibility follows from
uniqueness of the lifted curve.
\end{proof}

\subsubsection*{Candidate B.2e: finite quotients of a smooth model}
Let $k$ be algebraically closed, let $Y$ be a smooth projective integral
$k$-scheme, and let a finite group $\Gamma$ act faithfully on $Y$.
Set $K=k(Y)^\Gamma$.  The direct system of
$\operatorname{Pic}^{\tau}$ of proper normal models of $K/k$ is
essentially constant.
\begin{proof}
The quotient $T=Y/\Gamma$ is a normal projective scheme and
$q:Y\to T$ is finite.  We write $n=|\Gamma|$ and
$P=\operatorname{Pic}^{\tau}_{Y/k}$.
For a normal proper modification $W\to T$, let $Y_W$ be the
normalization of $W$ in $k(Y)$.  It is finite over $W$ and has a
proper birational morphism to $Y$; its $\Gamma$-quotient is $W$.
By B.1, pullback identifies $\operatorname{Pic}^{\tau}_{Y_W/k}$
with $P$, compatibly with $\Gamma$ and with further modifications.
Thus pullback gives a map
\[
 h_W:P_W:=\operatorname{Pic}^{\tau}_{W/k}\longrightarrow
 A:=P^\Gamma.
\]
Here $A$ is the closed fixed subgroup scheme.

The kernel of $h_W$ embeds in
$D(\Gamma^{\mathrm{ab}})=\underline{\operatorname{Hom}}(\Gamma,\mathbb G_m)$.
Indeed, after an fppf base change on a test scheme, a line bundle in
the kernel has a trivialization on $Y_W$.  Its canonical
$\Gamma$-linearization becomes a character, because
$\Gamma(Y_W\times_kS,\mathcal O)=\Gamma(S,\mathcal O)$ universally.
Changing the trivialization by a scalar does not change this character.
If it is trivial, the trivializing section and its inverse are invariant;
they descend because $(q_*\mathcal O_{Y_W})^\Gamma=\mathcal O_W$
universally after base change over $k$.  The descended section
trivializes the original line bundle.  This proves injectivity as
fppf sheaves.  The kernel is proper, so this monomorphism into the
finite group $D(\Gamma^{\mathrm{ab}})$ is a closed immersion.
Its order is at most $|\Gamma^{\mathrm{ab}}|\le n$.

Let $B_T\subset A$ be the image of $h_T$.  We claim that
$A_{\mathrm{red}}^0\subset B_T$.  Take $[L]\in A(k)$ and choose
a Cartier divisor $D$ representing $L$ on the regular scheme $Y$.
The divisor $D_\Sigma=\sum_{\gamma\in\Gamma}\gamma D$ is invariant.
At a prime divisor above a prime divisor of $T$, its coefficient is
divisible by the order of the decomposition group, hence by the
ramification index.  Therefore $D_\Sigma=q^*E$ for an integral Weil
divisor $E$ on $T$.  The divisor $nE$ is Cartier.  To verify this at
$x\in T$, trivialize $D_\Sigma$ on the semilocal scheme at the
finitely many points above $x$ and represent it by a rational
function $u$.  After shrinking around $x$, its norm has divisor
$q_*D_\Sigma=nE$.  Thus $nE$ is locally principal.
Since $[L]$ is $\Gamma$-invariant, we have
\[
 q^*\mathcal O_T(nE)\simeq L^{\otimes n^2}.
\]
If $L$ is torsion-equivalent to zero, this pullback is numerically
trivial.  Finite surjective pullback detects numerical triviality,
by degrees on integral curves.  Numerical triviality is equivalent
to torsion equivalence for projective schemes
\cite[Theorem 6.3]{Kleiman2005}.  Hence
$\mathcal O_T(nE)\in P_T(k)$, and $n^2[L]\in B_T(k)$.
Multiplication by $n^2$ is surjective on the abelian variety
$A_{\mathrm{red}}^0$, so the claim follows by density of its $k$-points.

The quotient $A/B_T$ is therefore a fixed finite group scheme, of
some order $N$.  Since $P_T\hookrightarrow P_W$ is a closed immersion,
the quotient $P_W/P_T$ maps to $A/B_T$; its kernel is a quotient
of $\ker(h_W)$.  Thus it is finite of order at most $nN$.
As $W$ varies, these finite quotients form a directed system of
closed immersions with uniformly bounded orders.  Choose one whose
order is maximal.  Every later closed immersion has equal order,
hence is an isomorphism.  This proves eventual constancy.  Normalized
graphs show that modifications of $T$ form a cofinal subsystem of
all proper normal models of $K/k$.
\end{proof}

\subsubsection*{Candidate B.2f: descent through a purely inseparable extension}
Let $k$ be algebraically closed of characteristic $p>0$, let $K'/K$
be a finite purely inseparable extension of finitely generated
$k$-fields, and assume that the Picard system for $K'/k$ is
essentially constant.  Then the same holds for $K/k$.
\begin{proof}
Choose $e$ with $(K')^{p^e}\subset K$, and put $q=p^e$.
Choose a normal projective model $T'$ of $K'$ after which its
Picard system is constant.  The relative Frobenius twist
$(T')^{(q)}$ has function field identified with $(K')^q$;
normalizing it in the intermediate field $K$ gives a normal
projective model $T$.  The normalization of $T$ in $K'$ is
exactly $T'$.  Thus $\pi:T'\to T$ is finite and purely inseparable.

For any normal proper modification $W\to T$, its normalization
$W'$ in $K'$ is finite over $W$ and maps birationally to $T'$.
The assumed stabilization identifies
$P_{W'}=\operatorname{Pic}^{\tau}_{W'/k}$ with the fixed group
$P'=\operatorname{Pic}^{\tau}_{T'/k}$.  Write
$h_W:P_W\to P'$ for pullback, and $B_T$ for the image of $h_T$.

The field inclusion $(K')^q\subset K$ and normality imply
$\mathcal O_{W'}^q\subset\mathcal O_W$ locally.  They give a
morphism $r_W:W\to(W')^{(q)}$ with
\[
 r_W\pi_W=F^e_{W'/k},\qquad
 \pi_W^{(q)}r_W=F^e_{W/k}.
\]
The two factorizations of relative Frobenius imply, on Picard
group schemes, that $[q]$ kills $\ker h_W$ and that
$[q]P'\subset B_T$.  More explicitly, relative Frobenius on a
scheme induces Verschiebung on its Picard scheme, and composing
this map with relative Frobenius of the Picard scheme is $[q]$;
the displayed factorizations give the two assertions by
functoriality.  Equivalently they follow directly by raising
transition functions of line bundles and their trivializations
to their $q$-th powers.

The second assertion contains $(P')_{\mathrm{red}}^0$ in $B_T$,
since $[q]$ is surjective on that abelian variety.  Therefore
$P'/B_T$ is finite, of some fixed order $N$.
Apply B.2d with $M=T$.  After the one fixed extension $E/k$,
every $P_W$ embeds in the same Jacobian $J$.  Since $[q]$
kills $\ker h_W$, its base change embeds in $J[q]$, whose
order is $q^{2g}$, where $g=\dim J$.  Thus
$|\ker h_W|\le q^{2g}$ uniformly in $W$.

The proper quotient $P_W/P_T$ maps to $P'/B_T$ with kernel a
quotient of $\ker h_W$, so it is finite of order at most
$q^{2g}N$.  The maximal-order argument of B.2e now proves
essential constancy over modifications of $T$, hence over all
normal proper models of $K/k$.
\end{proof}

\subsubsection*{Candidate B.2g: full stabilization over an algebraically closed field}
If $k$ is algebraically closed of positive characteristic and
$K/k$ is finitely generated, the full direct system
$\operatorname{Pic}^{\tau}_{X/k}$ over proper normal models of
$K/k$ is essentially constant.
\begin{proof}
By the regular Galois alteration theorem of de Jong
\cite[Situation 5.3, Theorem 5.13 and Corollary 5.15, pp.~613, 619--620]{deJong1997},
there are a smooth projective integral $Y$ and an action of a
finite group $\Gamma$ on $Y$ such that
$K'=k(Y)^\Gamma$ is finite purely inseparable over $K$.
The construction is projective: its semistable fibrations and
the equivariant modifications in Proposition 5.11 are projective,
starting with a projective model of $K$.
Dividing $\Gamma$ by the kernel of its action does not change
$K'$.  B.2e proves stabilization for $K'/k$, and B.2f descends
it to $K/k$.
\end{proof}
